% Solutions of Book 4, chapter 12 (caches).
\section*{Chapter \ref{ch:b2:caches} --- Caches and the Memory Hierarchy}

\begin{solution}{exo:b2:caches:1}
$128 / 16 = 8$ sets: 4 offset bits (bits 3 to 0), 3 index bits (bits 6 to 4), and the tag is the address shifted
right by 7. The trace:
\code{0x00} set 0, tag 0, \omterm{def:b2:caches:cache}{miss}; \code{0x04} and \code{0x08} \omterm{def:b2:caches:cache}{hit} in the same line; \code{0x10} set 1, tag 0, \omterm{def:b2:caches:cache}{miss};
\code{0x80} set 0, tag 1, \omterm{def:b2:caches:cache}{miss}, evicting tag 0; \code{0x00} \omterm{def:b2:caches:cache}{miss}, evicting tag 1; \code{0x14} set 1, tag 0, \omterm{def:b2:caches:cache}{hit};
\code{0x90} set 1, tag 1, \omterm{def:b2:caches:cache}{miss}; \code{0x1C} set 1, tag 0, \omterm{def:b2:caches:cache}{miss}; \code{0x84} set 0, tag 1, \omterm{def:b2:caches:cache}{miss}. Three \omterm{def:b2:caches:cache}{hits} and seven
misses, \code{MHHMMMHMMM}, as the chapter's tests check. The \omterm{def:b2:caches:three-c}{compulsory} misses are the first accesses to the four
lines \code{0x00}, \code{0x10}, \code{0x80} and \code{0x90}; the three others are conflicts, since only four
distinct lines are used and the \omterm{def:b2:caches:cache}{cache} has eight.
\end{solution}

\begin{solution}{exo:b2:caches:2}
The 512-byte array is 16 lines. With the correct mapping they go to 16 different sets of the 32: the first reading
misses once per line (16 misses, 112 \omterm{def:b2:caches:cache}{hits}) and the second \omterm{def:b2:caches:cache}{hits} everywhere. With the wrong mapping the index is the
top 5 bits of the address, the same, 2, for every address of the array: all 16 lines compete for one set of one
line. Each line still serves its 8 words in a row, so each reading misses 16 times: 32 misses. Taking the index from
the low bits above the offset spreads consecutive lines over all the sets, which is exactly what \omterm{def:b2:caches:locality}{spatial locality}
needs; the high bits change only between distant regions of memory.
\end{solution}

\begin{solution}{exo:b2:caches:3}
\omcode{code/bachelor-2/12-caches/exercises/flush.py}{4}{11}{Writes still owed by a \omterm{def:b2:caches:write}{write-back} \omterm{def:b2:caches:cache}{cache}}
Ten stores to one word in a \omterm{def:b2:caches:write}{write-back}, \omterm{def:b2:caches:write}{write-allocate} \omterm{def:b2:caches:cache}{cache}: the first misses and brings the line, marked dirty;
the nine others \omterm{def:b2:caches:cache}{hit}. Nothing is written to memory until the line is evicted or the \omterm{def:b2:caches:cache}{cache} emptied: one write in all,
which \code{total_writes} counts. A \omterm{def:b2:caches:write}{write-through} \omterm{def:b2:caches:cache}{cache} writes memory ten times. The tests check both counts.
\end{solution}

\begin{solution}{exo:b2:caches:4}
$32\,768 / (8 \times 64) = 64$ sets: 6 offset bits, 6 index bits, $32 - 12 = 20$ tag bits. The L2 of
\cref{out:b2:caches:getconf} has $2\,097\,152 / (16 \times 64) = 2048$ sets: 6 offset bits, 11 index bits, and with
32-bit addresses $32 - 17 = 15$ tag bits per line (more on the laptop itself, whose addresses are wider).
\end{solution}

\begin{solution}{exo:b2:caches:5}
8 sets, 4 offset bits, index bits 6 to 4, tag the address shifted right by 7. Every address of the trace has index 0,
\code{0x084} included ($\code{0x84} = 1000\,0100_2$), with tags 0, 2, 0, 4, 2, 1, 0, 6, 0. Set 0, \omterm{def:b2:caches:lru}{least recently
used} first: $[0]$, $[0, 2]$, $[2, 0]$ (\omterm{def:b2:caches:cache}{hit}), $[0, 4]$, $[4, 2]$, $[2, 1]$, $[1, 0]$, $[0, 6]$, $[6, 0]$ (\omterm{def:b2:caches:cache}{hit}): two \omterm{def:b2:caches:cache}{hits}.
In the \omterm{def:b2:caches:mapping}{direct-mapped cache} of 256 bytes (16 sets, index bits 7 to 4), \code{0x000}, \code{0x100}, \code{0x200} and
\code{0x300} all have index 0 and \code{0x084} index 8: every access evicts the one before in set 0, and all nine
\omterm{def:b2:caches:cache}{miss}. The 2-way \omterm{def:b2:caches:cache}{cache} does better because the trace alternates between two lines of the same set, which it can hold
together.
\end{solution}

\begin{solution}{exo:b2:caches:6}
\omcode{code/bachelor-2/12-caches/listings/exo_colsums.c}{26}{33}{Column sums, row by row}
\omoutput{code/bachelor-2/12-caches/outputs/exo_colsums.txt}{One run on the author's laptop, $8192 \times 4096$
grid.}
The slow version walks a column, 16 KiB from one element to the next: by \cref{prop:b2:caches:traversal}, with 8192
rows, every load misses in L1. The fast one walks the grid in the order it is stored, and the 4096 running sums,
32 KiB, stay in the L1 between rows. The same sums come out about thirty times faster.
\end{solution}

\begin{solution}{exo:b2:caches:7}
\omcode{code/bachelor-2/12-caches/listings/three_c.py}{5}{23}{Classifying misses with a shadow \omterm{def:b2:caches:cache}{cache}}
A line seen for the first time is \omterm{def:b2:caches:three-c}{compulsory}, whatever the shadow did. Otherwise, if the shadow, \omterm{def:b2:caches:mapping}{fully associative}
with the same number of lines, also missed, the \omterm{def:b2:caches:cache}{miss} is one of capacity; if it \omterm{def:b2:caches:cache}{hit}, the mapping alone caused the
\omterm{def:b2:caches:cache}{miss}. Each \omterm{def:b2:caches:cache}{miss} goes to exactly one branch, so the counts add up to the misses; the tests check it on random traces
for four associativities, and \cref{out:b2:caches:three-c} uses it.
\end{solution}

\begin{solution}{exo:b2:caches:8}
First, a \omterm{def:b2:caches:mapping}{fully associative} LRU \omterm{def:b2:caches:cache}{cache} of $C$ lines always holds the (at most $C$) distinct lines most recently used.
By induction on the accesses: a \omterm{def:b2:caches:cache}{hit} only reorders them; a \omterm{def:b2:caches:cache}{miss} inserts the new line, which becomes the most recent,
and if the \omterm{def:b2:caches:cache}{cache} was full, evicts the \omterm{def:b2:caches:lru}{least recently used} of the $C$ it held, which is the one that is no longer
among the $C$ most recent.
Now let access $t$ be to $\ell_{t \bmod (C+1)}$. For $t \le C$, the line is used for the first time: a \omterm{def:b2:caches:cache}{miss}. For
$t > C$, the accesses $t - C, \dots, t - 1$ are to $C$ distinct lines, all different from $\ell_{t \bmod (C+1)}$,
since any $C + 1$ consecutive accesses of the sweep are to distinct lines. The \omterm{def:b2:caches:cache}{cache} holds exactly those $C$ lines,
so access $t$ misses.
Over $C + 2$ lines, or any number more than $C$, the same argument applies: the last $C$ accesses are to $C$ distinct
lines other than the requested one. Over $C$ lines, the first $C$ accesses \omterm{def:b2:caches:cache}{miss} and all others \omterm{def:b2:caches:cache}{hit}, since the whole
sweep fits.
\end{solution}

\begin{solution}{exo:b2:caches:9}
\omcode{code/bachelor-2/12-caches/listings/transpose.c}{23}{30}{The tiled transpose}
\omoutput{code/bachelor-2/12-caches/outputs/transpose.txt}{One run on the author's laptop, $8192 \times 8192$
\code{int}.}
The plain version reads $a$ by rows but writes $b$ by columns, 32 KiB apart: every write is to a new line, and with
8192 rows, each line of $b$ is evicted before the next element of it is written, by
\cref{prop:b2:caches:traversal}. Exchanging the loops only moves the problem to the reads. A tile of $t \times t$ uses
$t$ lines of $a$ and $t$ lines of $b$, each reused $t$ times while it stays in \omterm{def:b2:caches:cache}{cache}; tiles of 8 make the transpose
about seven times faster here. A likely reason why larger tiles do worse: rows 32 KiB apart are a multiple of the L1's 4 KiB of sets
($64 \times 64$ bytes), so all the lines of a tile's column of $b$ fall in the same set of the L1, and 12 ways hold
the 8 lines of a tile of 8 but not the 16 of a tile of 16. Those tiles still gain from the L2.
\end{solution}

\begin{solution}{exo:b2:caches:10}
A column steps by $4N$ bytes, that is by $s = 4N / B$ lines, and the line of row $i$ has index $(i s + c) \bmod C$
for a constant $c$. If $s \ge C$, every row of the column falls in the same set; otherwise the column visits $C / s$
sets in turn (all powers of two). So a column uses $u = \min(N, \max(1, C B / (4N)))$ sets.
The line of $(i, j)$ is next needed at $(i, j + 1)$, which lies in the same line when $B \ge 8$. Between the two
accesses, the column order reads rows $i + 1, \dots, N - 1$ of column $j$ and rows $0, \dots, i - 1$ of column
$j + 1$. If $u < N$, rows $i + u$ or $i - u$ (one of them exists, since $u \le N / 2$) share the set of row $i$, and
their line, read in between, evicts it: every read misses. If $u = N$, every row has its own set, and the column order
misses only on the $4N^2 / B$ first uses. Now $u = N$ means $C B \ge 4N^2$: the \omterm{def:b2:caches:cache}{cache} must be as large as the whole
matrix. For the chapter's example, $C B = 1$ KiB and $4N^2 = 4$ KiB, so every read misses. A \omterm{def:b2:caches:mapping}{fully associative cache}
needs only $C \ge N$ lines (\cref{prop:b2:caches:traversal}), one column's worth: the direct mapping costs a factor
$4N / B$ in size on this access pattern.
\end{solution}

\begin{solution}{exo:b2:caches:11}
\omoutput{code/bachelor-2/12-caches/outputs/run_block_3c.txt}{The same misses, split by the classifier of
\cref{exo:b2:caches:7}.}
The \omterm{def:b2:caches:cache}{cache} has 64 lines in 16 sets; a row of either matrix is 512 bytes, 8 lines, so going down a column adds 8 to
the index modulo 16: a column uses only 2 sets, 8 lines.
Tile 1 is the plain transpose: $a$ is read by rows (1024 misses, one per line) and $b$ written by columns, each
column touching 128 lines, more than the \omterm{def:b2:caches:cache}{cache} holds, so every write misses: $1024 + 16\,384 = 17\,408$. With tiles of
$t \le 8$, the writes into one line of $b$ come $t$ at a time, and the misses of $b$ fall to about $16\,384 / t$:
$9216 = 1024 + 8192$ for $t = 2$. For $t = 4$ and 8, the lines of a tile of $a$ and of the matching tile of $b$ fall in
the same sets for some tiles, and the classifier finds 288, then 1824 conflict misses. At $t = 16$ a tile's column of $b$
has 16 lines in 2 sets of 4 ways: each line is evicted before its next write, and everything misses again, now as
conflicts, since a \omterm{def:b2:caches:mapping}{fully associative cache} would hold the tile. The least possible is the \omterm{def:b2:caches:three-c}{compulsory} count, every
line of both matrices once: $2 \times 65\,536 / 64 = 2048$, $0.125$ per element. Tile 8 is the best here, at 4896.
\end{solution}

\begin{solution}{exo:b2:caches:12}
Without L2, $\mathrm{AMAT} = 1 + 0.05 \times 100 = 6$\,ns. With an L2 of local \omterm{def:b2:caches:cache}{miss} ratio $m_2$,
$\mathrm{AMAT} = 1 + 0.05\,(4 + 100\,m_2)$, lower than 6 if and only if $4 + 100\,m_2 < 100$, that is
$m_2 < 0.96$: the L2 pays for itself as soon as it catches more than 4\,\% of the L1's misses. With $m_2 = 0.25$,
$\mathrm{AMAT} = 1 + 0.05 \times 29 = 2.45$\,ns. Adding an L3 replaces the term $100\,m_2$ by $m_2(15 + 100\,m_3)$,
lower if and only if $m_3 < 0.85$. With $m_3 = 0.4$:
$\mathrm{AMAT} = 1 + 0.05\,(4 + 0.25\,(15 + 40)) = 1 + 0.05 \times 17.75 = 1.8875$\,ns. The three AMATs are 6, 2.45
and about 1.89\,ns: each level gains less than the one before, because it only acts on what the levels above it \omterm{def:b2:caches:cache}{miss}.
\end{solution}

\begin{solution}{pb:b2:caches:1}
\textbf{1.} Nine reads per output pixel, and $(N - 2)^2$ output pixels: $9(N - 2)^2$, that is
$9 \times 254^2 = 580\,644$ for $N = 256$.

\textbf{2.} Rows $y - 1$, $y$ and $y + 1$. Row $y - 1$ read rows $y - 2$, $y - 1$ and $y$: two of the three were
already read.

\textbf{3.} A line of input row $r$ is read for output rows $r - 1$, $r$ and $r + 1$, three consecutive rows. Between
two consecutive uses, at output pixel $x$ of one row and output pixel $x$ of the next, the blur reads, at each column
position, three lines: those of three rows to the right of $x$ in the first output row and to its left in the next.
So fewer than three rows' worth of distinct lines, plus a few at the ends, come in between. If the \omterm{def:b2:caches:cache}{cache} holds three
rows, LRU keeps the line, and only its first use misses.

\textbf{4.} Every line misses once, so there are $N^2 / 16$ misses for $N^2$ pixels of the image: one every 16
pixels, a \omterm{def:b2:caches:cache}{miss} ratio of $1/16$ per pixel; since each pixel is read about nine times, $1/144$ per read.

\textbf{5.} A row is $256 \times 4 = 1024$ bytes, 16 lines; three rows are 48 lines, less than the 64 of a 4 KiB
\omterm{def:b2:caches:cache}{cache}. The 16 lines of a row fall in the 16 sets, one each, so three rows put 3 lines in each 4-way set: they fit.

\textbf{6.} One column of output reads three columns of input over all $N$ rows: at least one line per row, 256
lines, and two per row where $x - 1$ or $x + 1$ is in another line. The \omterm{def:b2:caches:cache}{cache} has 64 lines: the lines of the first
rows are gone when the next column needs them.

\textbf{7.} $9 \times 254^2 = 580\,644$, the simulator's figure.

\textbf{8.} The image is $256^2 \times 4 / 64 = 4096$ lines, and the row order misses once on each, by question 3.
The last column divides by the $254^2$ output pixels, not by the $256^2$ pixels read: $4096 / 64\,516 = 0.0635$,
against $4096 / 65\,536 = 1/16$.

\textbf{9.} Output pixel $(y, x)$ comes after $(y - 1, x)$, which read rows $y - 2$ to $y$; the lines of rows
$y - 1$ and $y$ are still in \omterm{def:b2:caches:cache}{cache}, and those of row $y + 1$ are new. That is one line if $x - 1$, $x$ and $x + 1$ are
in the same line, that is $x \bmod 16 \notin \{0, 15\}$, and two otherwise: on average $1 + 2/16 = 1.125$ misses per
pixel. Exactly: each of the 224 columns of the first kind misses $3 + 253 = 256$ times (three rows at $y = 1$, then
one per row), each of the 30 others 512 times: $224 \times 256 + 30 \times 512 = 72\,704$, the simulator's count, and
$72\,704 / 64\,516 = 1.1269$.

\textbf{10.} 4096 of them are \omterm{def:b2:caches:three-c}{compulsory}, the first use of each line. The others are capacity misses: the line of
$(y, x)$ is next needed one column later, after the lines of about 255 other rows, and a \omterm{def:b2:caches:mapping}{fully associative cache} of
64 lines would not keep it either.

\textbf{11.} No: by \cref{prop:b2:caches:traversal}, a column needs $C \ge N$, here 256 lines, and the \omterm{def:b2:caches:cache}{cache} has 64.

\textbf{12.} 4096: every line must be read at least once, and the row order reaches that bound.

\textbf{13.} Three: the strip's own line, and the lines holding $x_0 - 1$ and $x_0 + 16$.

\textbf{14.} $3N$ misses per strip, $N / 16$ strips, less the missing left neighbour of the first strip and right
neighbour of the last, $N$ lines each: $3N^2/16 - 2N = 12\,288 - 512 = 11\,776$.

\textbf{15.} Six lines per row of a strip of 64 (four of its own, two neighbours): $6N \times N/64 - 2N =
6144 - 512 = 5632$, and $5632 / 64\,516 = 0.0873$.

\textbf{16.} $(w/16 + 2) N$ misses per strip of $w N$ pixels, so $(w/16 + 2)/w = 1/16 + 2/w$ per pixel, which tends to
$1/16$, the row order's ratio, as $w$ grows.

\textbf{17.} Three rows of the strip need $3(w/16 + 2)$ lines; with 64 lines, $w/16 + 2 \le 21$, so $w \le 304$.
But rows 16 lines apart start in the same set: the $w/16 + 2$ consecutive lines of a row occupy as many consecutive
sets, and three rows put three lines in each of them only if a row covers at most 16 sets. A row of more than 16
lines puts 6 lines in some 4-way sets. The widest strip is therefore $w/16 + 2 = 16$, $w = 224$ pixels.

\textbf{18.} Each strip also loads the lines of its two neighbours, which the neighbouring strips load again. A strip
of 16 pays two extra lines for one useful one, $3/16$ per pixel; a strip of 64 two for four, $6/64$.

\textbf{19.} $4096 \times 4 = 16$ KiB per row; three rows are 48 KiB, exactly the L1, which also holds the output
and the stack, so the L1 is just short. A column's lines, 4096 of 64 bytes and more where neighbours cross a line,
are 256 to 512 KiB: far more than the L1, much less than the 2 MiB L2.

\textbf{20.} $72\,704 / 4096 \approx 17.8$ times more misses; measured, $483.7 / 53.7 \approx 9.0$ times slower.

\textbf{21.} The simulator has one \omterm{def:b2:caches:cache}{cache} of 4 KiB and counts every \omterm{def:b2:caches:cache}{miss} as a trip to memory. On the laptop the column
order's lines fit in the L2: its misses cost the L2's few nanoseconds, not a trip to memory. Both orders also pay the
same arithmetic, which dilutes the ratio, and the processor overlaps independent misses.

\textbf{22.} Measured: rows 53.7, strips of 64 69.3, strips of 16 170.9, columns 483.7\,ms. Predicted per pixel:
$0.0635 < 0.0873 < 0.1825 < 1.1269$. The order is the same.

\textbf{23.} $4094^2 \approx 16.8$ million pixels: about 3.2\,ns per pixel in row order, 28.9\,ns in column order, a
difference of about 25.7\,ns. The column order also writes its output down a column, so it misses about $1.125$ times
per pixel on the reads and once on the writes: about 12\,ns per \omterm{def:b2:caches:cache}{miss}, between the L2's and the L3's plateaus of
\cref{out:b2:caches:sweep}. It is an estimate: overlap and prefetching blur it.

\textbf{24.} In row order the blur misses once every 16 pixels, a predicted \omterm{def:b2:caches:cache}{miss} ratio of $1/16$, the least any order
can reach; the column order misses more than once per pixel ($1.127$ in the simulator, about eighteen times more),
strips of $w$ columns sit between at $1/16 + 2/w$; and on the laptop the row order is measured nine times faster than
the column order, with the four measured times in the order of the predicted \omterm{def:b2:caches:cache}{miss} ratios.
\end{solution}
